\documentclass[11pt]{article}
\usepackage{amsmath,amssymb,amsthm}

\newtheorem{lemma}{Lemma}

\begin{document}

\begin{lemma}[Unique Extension of Rooted Tree Isomorphism]
Let $(T_1,r_1),(T_2,r_2)$ be rooted trees of maximum degree $\Delta$.
Let $\varphi_0:\{r_1\}\to\{r_2\}$ be the root map.
Assume for every vertex $v\in T_1$ the multiset of isomorphism types of rooted subtrees at children of $v$
coincides with that of $\varphi(v)$ in $T_2$.
Then there exists a unique isomorphism $\varphi:T_1\to T_2$ extending $\varphi_0$.
\end{lemma}

\begin{proof}
Define $\varphi$ inductively by depth. At depth $0$, $\varphi(r_1)=r_2$.
Assume $\varphi$ defined on all vertices at depth $\le d$.
Let $v$ be at depth $d$, $\varphi(v)=w$.
Let $C(v),C(w)$ be children sets. By hypothesis, there exists a bijection
$\psi_v:C(v)\to C(w)$ preserving rooted subtree types.
For $u\in C(v)$, define $\varphi(u):=\psi_v(u)$.
This defines $\varphi$ on depth $d+1$.
By induction, $\varphi$ exists on all vertices.
Uniqueness follows since $\psi_v$ is uniquely determined at each step.
\end{proof}

\begin{lemma}[Transitivity of CFI Gadget Automorphisms]
Let $G_{\mathrm{CFI}}(v)$ be the CFI gadget at vertex $v$ of degree $d$.
Let $S_v$ be the set of local states (valid parity assignments).
Then $\mathrm{Aut}(G_{\mathrm{CFI}}(v))$ acts transitively on $S_v$.
\end{lemma}

\begin{proof}
Model $S_v=\{x\in\mathbb{F}_2^d:\sum x_i=0\}$.
For $i\neq j$, define $\tau_{ij}$ by $\tau_{ij}(x)=x+e_i+e_j$.
These preserve parity, hence lie in $\mathrm{Aut}$.
They generate translations $x\mapsto x+y$ for all $y\in S_v$.
Thus any $x'$ equals $x+y$ for some $y$, proving transitivity.
\end{proof}

\end{document}
